\chapter{Обзор существующих подходов к синхронизации состояния}

Механизмы синхронизации состояния определяют, каким образом сервер и клиенты обеспечивают
актуальность игрового мира, минимизируя расхождения между расчётами на стороне сервера и
восприятием игрока. Ниже рассмотрены основные методы,
используемые в современных сетевых играх различного жанра.

\section{Периодическая репликация состояния (State Snapshots)}

При использовании снимков сервер с фиксированной частотой отправляет клиентам актуальное состояние
мира или его части\cite{valvesource}. Снимок включает:

\begin{itemize}
    \item координаты и характеристики активных объектов;
    \item состояние игровых сущностей;
    \item флаги логики.
\end{itemize}

Этот подход используется в шутерах и королевских битвах\cite{unity_netcode_docs}, где важна высокая частота обновления положения объектов и событий.


Преимущества:

\begin{itemize}
    \item простота реализации;
    \item высокая точность согласования.
\end{itemize}

Недостатки:

\begin{itemize}
    \item высокий объем трафика при большом числе объектов;
    \item возможные скачки состояния при обновлениях;
    \item чувствительность к потере последовательных пакетов.
\end{itemize}

\section{Репликация только изменений (Delta Compression)}

При этом методе сервер отправляет не полный снимок, а только параметры, которые изменились со времени
последнего известного клиенту состояния\cite{unreal_replicate}.

Данный метод применяется в MMO и Survival-играх, где полный снимок мира не нужен.


Преимущества:

\begin{itemize}
    \item снижение объёма передаваемых данных;
    \item эффективная работа в больших сценах.
\end{itemize}

Недостатки:

\begin{itemize}
    \item клиент должен хранить предыдущее состояние;
    \item возможное накопление ошибок при потере пакета;
    \item необходимость периодической отправки полного снимка для восстановления согласованности.
\end{itemize}

\section{Событийная модель (Event-Based Sync)}

В рамках событийной модели сервер отправляет клиентам сведения не о состоянии объектов,
а о произошедших событиях\cite{GameEventSync2008}.

Событийная модель чаще всего используется в RPG, MOBA и стратегиях, где основная часть взаимодействий описывается дискретными событиями.


Преимущества:

\begin{itemize}
    \item малый объем передаваемых данных;
    \item высокая эффективность для логических и дискретных действий;
    \item высокая совместимость с детерминированной симуляцией.
\end{itemize}

Недостатки:

\begin{itemize}
    \item необходимость строгого детерминизма клиента;
    \item трудности коррекции накопившихся ошибок;
    \item расхождения при потерях событий.
\end{itemize}

\section{Синхронизация по входным данным (Input Sync)}

Сервер рассылает клиентам не состояние, а входные команды игроков. Детерминированная симуляция\cite{photon_sync}
на всех системах должна приводить к одинаковому результату.

Этот подход наиболее характерен для стратегий в реальном времени (RTS), файтингов и симуляторов, где критично поддерживать строгий детерминизм симуляции.


Преимущества:

\begin{itemize}
    \item минимальное сетевое потребление;
    \item высокая согласованность при строгом детерминизме.
\end{itemize}

Недостатки:

\begin{itemize}
    \item сильная зависимость от производительности клиента;
    \item сложная реализация физики;
    \item необходимость синхронизации тиков и входных очередей.
\end{itemize}

\section{Предсказание состояния (Client-Side Prediction)}

Клиент предсказывает поведение объекта немедленно после ввода пользователя, не ожидая
ответа сервера\cite{claypool2010latency}. После получения серверных данных выполняется корректировка.

Механизм предсказания широко используется в динамичных шутерах и экшен-играх, требующих минимальной задержки отклика на ввод игрока.


Преимущества:

\begin{itemize}
    \item высокая отзывчивость управления;
    \item минимизация задержки для игрока.
\end{itemize}

Недостатки:

\begin{itemize}
    \item возможные рывки (snap-back) при корректировке;
    \item риск накопления ошибок при нестабильной сети.
\end{itemize}

\section{Откат сетевого кода (Rollback Netcode)}

Метод используется в играх с высокой чувствительностью к синхронизации действий. Клиент
предсказывает действия другого клиента, а после получения реальных данных сервер выполняет откат\cite{ggpo}
состояния и пересчёт симуляции.

Rollback используется в аркадных PvP-играх и проектах с высокой чувствительностью к точности столкновений.


Преимущества:

\begin{itemize}
    \item высокая точность взаимодействий в шутерах;
    \item минимизация влияния сетевой задержки.
\end{itemize}

Недостатки:

\begin{itemize}
    \item повышенная нагрузка на вычислительные ресурсы;
    \item возможные визуальные искажения.
\end{itemize}
